\section{Задание}

В рамках выполнения лабораторной работы \textit{по варианту №6} необходимо выполнить следующие задачи:

\begin{enumerate}
  \item Исследуйте функцию нескольких переменных $f(x, y) = 10^{-2}(8x^2 + 2xy + 43x + 10y + 15)$ в заданной области аналитически и программно:
  \begin{itemize}
    \item определите все седловые точки, проведите исследование на локальные экстремумы \textit{(укажите все или докажите, что их нет)};
    \item найдите глобальный минимум функции;
    \item напишите программный модуль на языке Python для построения линий уровня этой функции;
    \item сделайте выводы, какие сложности для оптимизационного алгоритма могут возникнуть в случае такой функции.
  \end{itemize}

  \item Реализуйте и исследуйте алгоритмы градиентного спуска и эвристику \textbf{Adam} для задачи преодоления седловой точки.
  \begin{itemize}
    \item выберите начальное приближение, скорость обучения и критерий останова так, чтобы последовательные приближения стандартного градиентного спуска сходились к седловой точке за $\approx 100$ итераций, зафиксируйте эти параметры и постройте визуализацию линий уровня и траектории;
    \item реализуйте программно на языке Python эвристику по варианту --- метод \textbf{Adam} --- и с его помощью за то же количество итераций получите последовательность приближений, которая «преодолевает» седловую точку;
    \item подберите область отображения для эффективной иллюстрации диаграммы последовательных приближений и, не меняя её, найдите значения гиперпараметров для двух случаев:
    \begin{itemize}
      \item когда явно просматривается пилообразное движение в стороны от направления к глобальному минимуму;
      \item когда направление движения более чёткое и без колебаний.
    \end{itemize}
    \item используйте готовую реализацию алгоритма \textbf{Adam} \textit{(рекомендуется библиотека PyTorch)} для подтверждения успешного решения задачи «преодоления» седловой точки и постройте один вариант визуализации диаграммы приближений в той же области.
  \end{itemize}
\end{enumerate}